% 2010 年江苏卷第 7 题：算法流程图。
% 来源 PDF：普通高考/2010/2010江苏.pdf，SHA-256 e701df878d51a700720549274fe7969bd54342f0f58acb051c516fc502b58af7。
% 原页：第 1 页（A3 三栏）第 1 栏下部，书页码 605。
% 重绘依据：开始 → S←1 → n←1 → S←S+2^n → n←n+1 → 判断 S≥33；
%   否（右）→ 回到 S←S+2^n；是（下）→ 输出 S → 结束。
\begin{tikzpicture}[semithick, >=Stealth,
  every node/.style={inner sep=1.5pt, font=\small}]
  \node (start) [draw, rounded corners=4pt, minimum width=1.8cm, minimum height=0.55cm] at (0,0) {开始};
  \node (s1) [draw, minimum width=1.8cm, minimum height=0.6cm] at (0,-1.3) {$S\leftarrow1$};
  \node (n1) [draw, minimum width=1.8cm, minimum height=0.6cm] at (0,-2.65) {$n\leftarrow1$};
  \node (add) [draw, minimum width=3.0cm, minimum height=0.6cm] at (0,-4.15) {$S\leftarrow S+2^n$};
  \node (inc) [draw, minimum width=1.8cm, minimum height=0.6cm] at (0,-5.5) {$n\leftarrow n+1$};
  \node (test) [draw, diamond, aspect=2.2, minimum width=2.6cm, minimum height=1.1cm] at (0,-7.15) {$S\geqslant33$};
  \node (out) [draw, trapezium, trapezium left angle=72, trapezium right angle=108,
    minimum width=2.2cm, minimum height=0.55cm] at (0,-8.9) {输出 $S$};
  \node (end) [draw, rounded corners=4pt, minimum width=1.8cm, minimum height=0.55cm] at (0,-10.2) {结束};

  \draw[->] (start) -- (s1);
  \draw[->] (s1) -- (n1);
  \draw[->] (n1) -- (add);
  \draw[->] (add) -- (inc);
  \draw[->] (inc) -- (test);
  \draw[->] (test.south) -- node[right] {Y} (out.north);
  \draw[->] (out) -- (end);
  \draw[->] (test.east) -- node[above] {N} (4.2,-7.15) -- (4.2,-3.5) -- (0,-3.5) -- (add.north);
\end{tikzpicture}
